% Problem record op_fe2c638b34652717. This TeX file is derived from
% database/problems_json/op_fe2c638b34652717.json by scripts/sync-tex.mjs; edit the JSON record
% and rerun the script rather than editing this file.
\section{Polynomial Pauli periodicity at a fixed Clifford hierarchy level}
\label{sec:op_fe2c638b34652717}

\paragraph{Problem.}
Can the Pauli periodicity of gates in one fixed level of the qubit
Clifford hierarchy above the Clifford group grow at least polynomially with the
number of qubits?  For $n\geq1$, let $\mathcal{P}_n$ be the $n$-qubit Pauli
group, generated by $iI$ and the single-qubit operators $X_j,Y_j,Z_j$, and let
$\mathbb{T}:=\{e^{i\theta}:\theta\in\mathbb{R}\}$.  Define the hierarchy modulo
global phase by
\begin{equation}
  \mathcal{C}_1(n):=\mathbb{T}\mathcal{P}_n,
  \qquad
  \mathcal{C}_{k+1}(n):=
  \bigl\{U\in U(2^n):UPU^\dagger\in\mathcal{C}_k(n)
  \ \text{for every}\ P\in\mathcal{P}_n\bigr\},
  \qquad k\geq1,
  \label{eq:fe2c-hierarchy}
\end{equation}
so that $\mathcal{C}_2(n)$ is the Clifford group.  For a unitary
$U\in U(2^n)$, its projective Pauli periodicity is
\begin{equation}
  m(U):=\min\bigl\{m\in\mathbb{Z}_{\geq0}:U^{2^m}\in\mathbb{T}\mathcal{P}_n\bigr\},
  \label{eq:fe2c-periodicity}
\end{equation}
with $m(U):=\infty$ when no such $m$ exists.  Thus $m(U)$ counts repeated
squarings rather than the total number of powers, and scalar phases are
ignored, so that $m(e^{i\theta}U)=m(U)$.  For a level $k$ of
Eq.~\eqref{eq:fe2c-hierarchy}, set
\begin{equation}
  M_k(n):=\sup\bigl\{m(U):U\in\mathcal{C}_k(n),\ m(U)<\infty\bigr\}.
  \label{eq:fe2c-max}
\end{equation}
The question asks whether the quantity in Eq.~\eqref{eq:fe2c-max} satisfies
\begin{equation}
  \exists\,k\geq3\ \exists\,a,c>0\ \forall\,N\geq1\ \exists\,n\geq N:
  \qquad M_k(n)\geq a\,n^{c}.
  \label{eq:fe2c-target}
\end{equation}
A negative answer means that $M_k(n)=o(n^{c})$ as $n\to\infty$ for every fixed
$k\geq3$ and every $c>0$.

\subsection*{Status}

Unsolved

\subsection*{Source}

Xu and Wang pose the question explicitly as Open Problem~2 in Section~5:
for a Pauli-periodic gate in a fixed, sufficiently high level of the
hierarchy, can the largest Pauli periodicity grow as $\operatorname{poly}(n)$
instead of $\log n$? \sourcecite{ref:fe2c-xw26}{XW26}.
Eqs.~\eqref{eq:fe2c-periodicity} and~\eqref{eq:fe2c-target} fix a
phase-insensitive definition and a precise reading of polynomial growth; the
restriction $k\geq3$ excludes the Clifford level, where the answer is known.

\subsection*{Progress}

\begin{itemize}
  \item The classification of diagonal hierarchy gates by Cui, Gottesman, and
  Krishna places the one-qubit phase gates
  \begin{equation}
    R_k:=\operatorname{diag}\bigl(1,e^{i\pi/2^{k-1}}\bigr)\in\mathcal{C}_k(1),
    \qquad k\geq2.
    \label{eq:fe2c-phase}
  \end{equation}
  Since $R_k^{2^{k-2}}=\operatorname{diag}(1,i)$ is not a Pauli operator up to
  phase while $R_k^{2^{k-1}}=Z$, the gates in Eq.~\eqref{eq:fe2c-phase} have
  $m(R_k)=k-1$, and tensoring with identities gives $M_k(n)\geq k-1$ for all
  $n$.  Periodicity therefore grows without bound when the level grows, but this
  does not address growth in $n$ at a fixed level
  \sourcecite{ref:fe2c-cgk17}{CGK17}.

  \item Xu and Wang proved in Theorem~4 that every Clifford gate $U$ with
  $m(U)<\infty$ satisfies $m(U)\leq\lceil\log_2(2n)\rceil$, with equality
  attained for every $n>1$; Proposition~9 in Section~3.2 gives explicit
  saturating $S_X$--CNOT--$H$ strings for $n>2$, and for $n=1$ the phase gate
  $\operatorname{diag}(1,i)$ attains the bound.  Hence
  \begin{equation}
    M_2(n)=\lceil\log_2(2n)\rceil\qquad(n\geq1).
    \label{eq:fe2c-clifford}
  \end{equation}
  The proof passes to the binary symplectic matrix $F$ of $U$:
  $U^{2^t}\in\mathbb{T}\mathcal{P}_n$ exactly when $F^{2^t}=I$, a matrix over
  $\mathbb{F}_2$ of two-power order is unipotent, and $(F-I)^{2n}=0$ yields the
  bound; tightness uses regular unipotent elements of
  $\operatorname{Sp}(2n,\mathbb{F}_2)$.  Because $F$ determines $U$ only up to a
  global phase, Eq.~\eqref{eq:fe2c-clifford} refers to the phase-insensitive
  periodicity of Eq.~\eqref{eq:fe2c-periodicity}; if Pauli membership allowed
  only the phases $\pm1,\pm i$, the scalar Clifford gates $e^{i\pi/2^{r}}I$
  would have periodicity $r-1$ for every $r\geq1$
  \sourcecite{ref:fe2c-xw26}{XW26}.

  \item Since $\mathcal{C}_2(n)\subseteq\mathcal{C}_k(n)$ for $k\geq2$,
  Eqs.~\eqref{eq:fe2c-phase} and~\eqref{eq:fe2c-clifford} give the elementary
  lower bound
  \begin{equation}
    M_k(n)\geq\max\bigl\{k-1,\lceil\log_2(2n)\rceil\bigr\},
    \qquad k\geq2,
    \label{eq:fe2c-lower}
  \end{equation}
  which is logarithmic in $n$ at fixed $k$
  \sourcecite{ref:fe2c-cgk17}{CGK17}, \sourcecite{ref:fe2c-xw26}{XW26}.  In
  Section~5, Open Problem~2, Xu and Wang ask for the largest Pauli periodicity
  in a fixed level and whether, for a sufficiently high fixed level, it can be
  polynomial rather than logarithmic in $n$; Eq.~\eqref{eq:fe2c-target} isolates
  that alternative \sourcecite{ref:fe2c-xw26}{XW26}.

  \item Ordinary finite order is a different criterion.  For $U$ with
  $m(U)<\infty$, let $o(U):=\min\{r\geq1:U^r\in\mathbb{T}I\}$ be its projective
  order.  Every element of $\mathbb{T}\mathcal{P}_n$ squares to a scalar, so
  $U^{2^{m(U)+1}}\in\mathbb{T}I$ and $o(U)$ divides $2^{m(U)+1}$; if $m(U)\geq1$,
  then $o(U)\leq2^{m(U)-1}$ would give $U^{2^{m(U)-1}}\in\mathbb{T}I$, contrary to
  the minimality in Eq.~\eqref{eq:fe2c-periodicity}.  Hence
  \begin{equation}
    o(U)\ \text{is a power of two},
    \qquad
    2^{m(U)}\leq o(U)\leq2^{m(U)+1}.
    \label{eq:fe2c-order}
  \end{equation}
  Consequently, Clifford circuits realizing linear-feedback shift registers of
  period $2^n-1$, which Xu and Wang contrast with Open Problem~2 in Section~5,
  have $m(U)=\infty$ for $n\geq2$ and do not bear on
  Eq.~\eqref{eq:fe2c-target} \sourcecite{ref:fe2c-xw26}{XW26}.
\end{itemize}

\subsection*{References}

\begin{enumerate}[label={},leftmargin=4.8em,labelsep=0.5em]
  \footnotesize
  \item[\textup{[CGK17]}]\label{ref:fe2c-cgk17}
  S. X. Cui, D. Gottesman, and A. Krishna, ``Diagonal gates in the Clifford
  hierarchy,'' \emph{Physical Review A} \textbf{95}, 012329 (2017).
  \href{https://doi.org/10.1103/PhysRevA.95.012329}{doi:10.1103/PhysRevA.95.012329};
  \href{https://arxiv.org/abs/1608.06596}{arXiv:1608.06596}.

  \item[\textup{[XW26]}]\label{ref:fe2c-xw26}
  Y. Xu and X. Wang, ``Controlled jump in the Clifford hierarchy,'' arXiv
  preprint (2026).
  \href{https://doi.org/10.48550/arXiv.2602.22201}{doi:10.48550/arXiv.2602.22201};
  \href{https://arxiv.org/abs/2602.22201}{arXiv:2602.22201}.
\end{enumerate}

\subsection*{Comment}

The open alternative is whether some fixed level $k\geq3$ contains gates
whose Pauli periodicity is at least polynomial in $n$ for infinitely many $n$,
or whether $M_k(n)=o(n^c)$ for every fixed $k\geq3$ and every $c>0$.  At fixed
$k$, the available lower bound, Eq.~\eqref{eq:fe2c-lower}, is logarithmic.  The
exact value in Eq.~\eqref{eq:fe2c-clifford} rests on the group structure and
symplectic representation of the Clifford level, which are not available at
levels $k\geq3$, since these levels are not groups.  Letting the level grow with
$n$, as for the gates in Eq.~\eqref{eq:fe2c-phase}, or using large projective
order without the two-power structure of Eq.~\eqref{eq:fe2c-order}, does not
meet the quantifiers of Eq.~\eqref{eq:fe2c-target}.  Xu and Wang note that
polynomial periodicity at a fixed level could reduce the qubit cost of
controlled-jump constructions, although their exact level formula for
controlled gates is proved only for Clifford targets
\sourcecite{ref:fe2c-xw26}{XW26}.  The level reachable by controlling targets above the Clifford
level, Open Problem~1 of Xu and Wang, is the record on
\href{https://qiqc-op.com/problem/op_e432fc5ff73e63be/}{the highest Clifford-hierarchy level of controlled gates with
higher-level targets}.  Literature checked through 15 September
2026; no construction with polynomial periodicity at a fixed level and no upper
bound excluding it was found.

\subsection*{Field}

Quantum algorithm

\subsection*{Topic}

Clifford hierarchy

\subsection*{ID}

\texttt{op\_fe2c638b34652717}
